% Source: Weinberg GR, physical PDF pp. 505--515
%         (printed pp. 481--491; boundary pages are partial).
% Coverage: Section 15.3, The Matter-Dominated Era; equations
%           (15.3.1)--(15.3.37), including all unnumbered displays.
% Figures/tables/footnotes: Figure 15.1; reference notes 15--25; no tables.
% Status: source-reviewed and compile-clean.
% Uncertainties: the source's final approximation in (15.3.34), t/3,
%                conflicts with the preceding expression (which evaluates
%                to 3t); it is preserved and marked VERIFY SOURCE below.
% MODERNIZATION: R_{\mathrm{source}}(t) -> a(t), with modern H and q
% notation throughout. The source's nonstandard spelling is regularized to
% "deceleration." Curvature signs do not appear explicitly
% beyond k here; the Friedmann equations already use the binding convention
% fixed and audited in Section 15.1.

\subsection{The Matter-Dominated Era}\label{sec:15.3}

We have noted that the energy density of the \emph{known} forms of radiation
in the present universe is less than one-hundredth the density of rest mass.
According to Eqs.~\eqref{eq:15.1.22} and \eqref{eq:15.1.23}, the energy
density of rest mass scales as \(a^{-3}\), and the energy density of
radiation scales as \(a^{-4}\), so we may conclude with some confidence that
the expansion of the universe has been governed by its nonrelativistic matter
content at least since the time when \(a(t)\) was one-hundredth its present
value. This period certainly goes back long before the emission of any of the
light collected at Mt.\ Palomar, for the most distant galaxies and
quasi-stellar objects observed have red shifts \(z\) that are much less than
100, and in fact less than 3! The study of the empirical relations between red
shifts, luminosities, numbers, angular diameters, and so on, can therefore
reveal only the matter-dominated era of the history of the universe.

The dynamical equation governing the universe during this era is Einstein's
equation \eqref{eq:15.1.20}:
\begin{equation}
  \dot a^2+k=\frac{8\pi G}{3}\rho a^2,
  \tag{15.3.1}\label{eq:15.3.1}
\end{equation}
with \(\rho\) taking the form \eqref{eq:15.1.22} appropriate to a
matter-dominated universe:
\begin{equation}
  \frac{\rho}{\rho_0}
  =
  \left(\frac{a}{a_0}\right)^{-3}.
  \tag{15.3.2}\label{eq:15.3.2}
\end{equation}
It is convenient to make use of Eqs.~\eqref{eq:15.2.5} and
\eqref{eq:15.2.6} to write \(\rho_0\) and \(k/a_0^2\) in terms of \(q_0\)
and \(H_0\):
\[
  \frac{k}{a_0^2}=(2q_0-1)H_0^2,
  \qquad
  \frac{8\pi G\rho_0}{3}=2q_0H_0^2.
\]
Equations~\eqref{eq:15.3.1} and \eqref{eq:15.3.2} then give
\begin{equation}
  \left(\frac{\dot a}{a_0}\right)^2
  =
  H_0^2
  \left[
    1-2q_0+2q_0\left(\frac{a_0}{a}\right)
  \right].
  \tag{15.3.3}\label{eq:15.3.3}
\end{equation}
The solution may in general be expressed as a formula for \(t\) in terms of
\(a\):
\begin{equation}
  t
  =
  \frac1{H_0}
  \int_0^{a/a_0}
  \left[
    1-2q_0+\frac{2q_0}{x}
  \right]^{-1/2}
  \dd x,
  \tag{15.3.4}\label{eq:15.3.4}
\end{equation}
with \(t=0\) defined as the time when \(a\ll a_0\). In particular, the
present age of the universe is
\begin{equation}
  t_0
  =
  \frac1{H_0}
  \int_0^1
  \left[
    1-2q_0+\frac{2q_0}{x}
  \right]^{-1/2}
  \dd x.
  \tag{15.3.5}\label{eq:15.3.5}
\end{equation}
For any positive \(q_0\), the age of the universe must be less than the
Hubble time,
\begin{equation}
  t_0<\frac1{H_0},
  \tag{15.3.6}\label{eq:15.3.6}
\end{equation}
as already remarked in Section~\ref{sec:15.1}.

The behavior of the result \eqref{eq:15.3.4} may conveniently be discussed
under three special cases (see Figure~\ref{fig:15.1}):

\emph{(A) \(q_0>\tfrac12\)
\((k=+1,\ \rho_0>\rho_c)\).} It is convenient here to define a
\emph{development angle} \(\theta\) by
\begin{equation}
  1-\cos\theta
  =
  \left(\frac{2q_0-1}{q_0}\right)\frac{a(t)}{a_0}.
  \tag{15.3.7}\label{eq:15.3.7}
\end{equation}

% Source: Weinberg GR, physical PDF p. 506 (printed p. 482).
% Coverage: Figure 15.1, matter-dominated Friedmann solutions.
% Status: source-reviewed and compile-clean; redrawn in TikZ with the
%         binding scale factor a.
% Uncertainties: none.

\begingroup
\renewcommand{\thefigure}{15.1}
\begin{figure}[htbp]
  \centering
  \begin{tikzpicture}[x=0.88cm,y=0.88cm,>=Stealth]
    \draw[->] (0,0) -- (7.45,0) node[right] {\(t\)};
    \draw[->] (0,0) -- (0,4.85) node[above] {\(a\)};

    % Open, flat, and closed matter-dominated Friedmann solutions.
    \draw[thick]
      plot[smooth] coordinates
      {(0,0) (0.28,0.95) (0.62,1.75) (1.05,2.45)
       (1.48,3.12) (1.92,3.78) (2.35,4.36)};
    \node[anchor=west] at (2.08,4.27) {\(k=-1\)};

    \draw[thick]
      plot[smooth] coordinates
      {(0,0) (0.32,0.85) (0.78,1.62) (1.38,2.35)
       (2.05,3.03) (2.75,3.62) (3.30,4.05)};
    \node[anchor=west] at (3.17,4.08) {\(k=0\)};

    \draw[thick]
      (0,0)
      .. controls (0.58,1.86) and (1.82,2.55) .. (3.55,2.62)
      .. controls (5.25,2.57) and (6.68,1.42) .. (7.15,0);
    \node[anchor=west] at (5.92,2.52) {\(k=+1\)};

    % Values of q along the k=-1 branch.
    \foreach \x/\y/\lab/\dx/\dy in {
      0.00/0.00/{0.5}/0.18/0.18,
      0.43/1.30/{0.3}/-0.48/0.07,
      1.04/2.43/{0.25}/-0.55/0.08,
      1.72/3.50/{0.21}/-0.58/0.08
    }{
      \fill[white] (\x,\y) circle (2.2pt);
      \draw (\x,\y) circle (2.2pt);
      \node at ({\x+\dx},{\y+\dy}) {\(\lab\)};
    }

    % Values of q along the k=+1 branch.
    \foreach \x/\y/\lab/\dx/\dy in {
      0.37/0.84/{1}/0.22/-0.24,
      1.13/1.83/{2}/0.24/-0.25,
      2.38/2.45/{7.5}/0.00/-0.32,
      3.55/2.62/{\infty}/0.00/-0.32
    }{
      \fill[white] (\x,\y) circle (2.2pt);
      \draw (\x,\y) circle (2.2pt);
      \node at ({\x+\dx},{\y+\dy}) {\(\lab\)};
    }
  \end{tikzpicture}
  \caption{Solutions of Einstein's equations for a Robertson--Walker
  universe with curvature \(k=+1\), \(k=0\), and \(k=-1\). The numbers
  along the curves \(k=\pm1\) give the values of the deceleration parameter
  \(q\) at various epochs.}
  \label{fig:15.1}
\end{figure}
\endgroup


Then Eq.~\eqref{eq:15.3.4} gives
\begin{equation}
  H_0t
  =
  q_0(2q_0-1)^{-3/2}(\theta-\sin\theta).
  \tag{15.3.8}\label{eq:15.3.8}
\end{equation}
This is the equation of a cycloid; \(a(t)\) increases from zero at
\(\theta=0,\ t=0\); reaches a maximum at
\begin{equation}
  \theta_m=\pi,\qquad
  t_m=\frac{\pi q_0}{H_0(2q_0-1)^{3/2}},\qquad
  a(t_m)=\frac{2q_0a_0}{2q_0-1},
  \tag{15.3.9}\label{eq:15.3.9}
\end{equation}
and then returns to zero at \(\theta=2\pi,\ t=2t_m\). The present instant
is defined by setting \(a(t)\) equal to \(a_0\) in
Eq.~\eqref{eq:15.3.7}; the present value of the development angle
\(\theta\) is then given by
\begin{equation}
  \cos\theta_0=\frac1{q_0}-1,
  \tag{15.3.10}\label{eq:15.3.10}
\end{equation}
and so the age of the universe is
\begin{equation}
  t_0
  =
  H_0^{-1}q_0(2q_0-1)^{-3/2}
  \left[
    \cos^{-1}\left(\frac1{q_0}-1\right)
    -\frac1{q_0}(2q_0-1)^{1/2}
  \right].
  \tag{15.3.11}\label{eq:15.3.11}
\end{equation}
For example, if we believe that \(q_0\approx1\) and
\(H_0^{-1}=13\times10^9\) years, then Eq.~\eqref{eq:15.3.10} gives
\(\theta_0\approx\pi/2\), Eq.~\eqref{eq:15.3.11} gives the age of the
universe now as
\begin{equation}
  t_0
  \approx
  \left(\frac{\pi}{2}-1\right)H_0^{-1}
  \approx7.5\times10^9\ \text{years},
  \tag{15.3.12}\label{eq:15.3.12}
\end{equation}
and Eq.~\eqref{eq:15.3.9} shows that the universe will reach its maximum
radius \(a(t_m)\approx2a_0\) at a time
\begin{equation}
  t_m\approx\pi H_0^{-1}\approx40\times10^9\ \text{years}.
  \tag{15.3.13}\label{eq:15.3.13}
\end{equation}
The whole life cycle of the universe takes a time \(2t_m\), or about
\(80\times10^9\) years.

\emph{(B) \(q_0=\tfrac12\)
\((k=0,\ \rho_0=\rho_c)\).} Here Eq.~\eqref{eq:15.3.4} gives
\begin{equation}
  \frac{a(t)}{a_0}
  =
  \left(\frac{3H_0t}{2}\right)^{2/3},
  \tag{15.3.14}\label{eq:15.3.14}
\end{equation}
so \(a(t)\) increases without limit. The age of the universe is given for
\(H_0^{-1}\approx13\times10^9\) years as
\begin{equation}
  t_0=\frac23H_0^{-1}\approx9\times10^9\ \text{years}.
  \tag{15.3.15}\label{eq:15.3.15}
\end{equation}
This is known as the \emph{Einstein--de Sitter model}.

\emph{(C) \(0<q_0<\tfrac12\)
\((k=-1,\ \rho_0<\rho_c)\).} The results \eqref{eq:15.3.7} and
\eqref{eq:15.3.8} can be applied here, except that now the development angle
\(\theta\) is imaginary,
\[
  \theta=\ii\Psi,
\]
and so
\begin{equation}
  H_0t
  =
  q_0(1-2q_0)^{-3/2}(\sinh\Psi-\Psi),
  \tag{15.3.16}\label{eq:15.3.16}
\end{equation}
with \(\Psi\) given by
\begin{equation}
  \cosh\Psi-1
  =
  \frac{1-2q_0}{q_0}\frac{a(t)}{a_0}.
  \tag{15.3.17}\label{eq:15.3.17}
\end{equation}
Just as in case (B), the scale factor \(a(t)\) here increases without limit;
for \(t\to\infty\), we have
\begin{equation}
  \frac{a(t)}{a_0}
  \longrightarrow
  \frac12q_0(1-2q_0)^{-1}e^\Psi
  \longrightarrow
  (1-2q_0)^{1/2}H_0t.
  \tag{15.3.18}\label{eq:15.3.18}
\end{equation}
The present moment is defined by setting \(a(t)\) equal to \(a_0\) in
Eq.~\eqref{eq:15.3.17}:
\begin{equation}
  \cosh\Psi_0=\frac1{q_0}-1,
  \tag{15.3.19}\label{eq:15.3.19}
\end{equation}
and the age of the universe is
\begin{equation}
  t_0
  =
  H_0^{-1}
  \left[
    (1-2q_0)^{-1}
    -q_0(1-2q_0)^{-3/2}
      \cosh^{-1}\left(\frac1{q_0}-1\right)
  \right].
  \tag{15.3.20}\label{eq:15.3.20}
\end{equation}
For instance, if we take the mass density of the universe to be that contained
within the galaxies, then according to Eq.~\eqref{eq:15.2.15},
\(q_0\) is about 0.014, so \(\Psi_0\approx5\), and the age of the universe
is nearly equal to the Hubble time:
\begin{equation}
  t_0\approx0.96H_0^{-1}\approx13\times10^9\ \text{years}.
  \tag{15.3.21}\label{eq:15.3.21}
\end{equation}

It is worth mentioning here that the deceleration parameter
\(q=-a\ddot a/\dot a^2\) will in general change with time. For \(k=+1\),
\(q(t)\) is given by the analogue of Eq.~\eqref{eq:15.3.10},
\[
  q=(1+\cos\theta)^{-1},
\]
so as \(\theta\) goes from 0 to \(2\pi\) during one cosmic cycle, \(q\)
rises from \(\tfrac12\) to infinity and then drops again to \(\tfrac12\).
For \(k=-1\), \(q(t)\) is given by the analogue of
Eq.~\eqref{eq:15.3.19}:
\[
  q=(1+\cosh\Psi)^{-1},
\]
so as \(\Psi\) goes from 0 to infinity, \(q\) drops steadily from
\(\tfrac12\) to 0. Only in the case \(k=0\) does \(q\) remain constant,
at the value \(q=\tfrac12\). Thus no special significance can be attached
to any particular value of \(q_0\) other than \(q_0=\tfrac12\). For
instance, the straight-line fit of luminosity distance versus red shift
indicates that \(q_0\approx1\) (unless evolutionary or selection effects are
important; see Section~\ref{sec:14.6}), but in a matter-dominated universe
this must be an accident, for if at present \(q_0=1\), then previously
\(q_0<1\), and in the future \(q_0>1\). It is only in a
radiation-dominated universe that \(k=0\) entails \(q_0=1\) [see
Eq.~\eqref{eq:15.2.17}], so that a deceleration parameter of unity is stable.

The formulas for \(a(t)\) derived above may be used to extend the
phenomenological analysis of the last chapter out to arbitrarily large red
shifts. According to Eq.~\eqref{eq:14.3.6}, light, which arrives at time
\(t_0\) with red shift \(z\), was emitted when the scale factor had the value
\begin{equation}
  a_1=\frac{a_0}{1+z}.
  \tag{15.3.22}\label{eq:15.3.22}
\end{equation}
The comoving radial coordinate of the light source is given by
Eqs.~\eqref{eq:14.3.1}, \eqref{eq:14.3.2}, and \eqref{eq:15.3.3}:
\begin{align*}
  \int_0^{r_1}\frac{\dd r}{\sqrt{1-kr^2}}
  &=
  \int_{t_1}^{t_0}\frac{\dd t}{a(t)}
  =
  \int_{a_1}^{a_0}\frac{\dd a}{a\dot a}\\
  &=
  \frac1{a_0H_0}
  \int_{(1+z)^{-1}}^1
  \left[
    1-2q_0+\frac{2q_0}{x}
  \right]^{-1/2}
  x^{-1}\dd x.
\end{align*}
With the aid of Eq.~\eqref{eq:15.2.5}, it is straightforward to show that
for all three possible values of \(k\), the formula for \(r_1\) is the same:
\begin{equation}
  r_1
  =
  \frac{
    zq_0+(q_0-1)\bigl(-1+\sqrt{2q_0z+1}\bigr)
  }{
    H_0a_0q_0^2(1+z)
  }.
  \tag{15.3.23}\label{eq:15.3.23}
\end{equation}
The ``luminosity distance,'' measured by comparison of apparent with absolute
luminosity, is then given by Eq.~\eqref{eq:14.4.14} as
\begin{equation}
  d_L
  =
  a_0r_1(1+z)
  =
  \frac1{H_0q_0^2}
  \left[
    zq_0+(q_0-1)\bigl(-1+\sqrt{2q_0z+1}\bigr)
  \right].
  \tag{15.3.24}\label{eq:15.3.24}
\end{equation}
In some determinations of \(q_0\), the observed \(d_L\) versus \(z\) curve
is compared with this exact formula, rather than the model-independent
approximation \eqref{eq:14.6.8},
\begin{equation}
  d_L
  \approx
  H_0^{-1}\left[z+\frac12(1-q_0)z^2\right],
  \tag{15.3.25}\label{eq:15.3.25}
\end{equation}
which is strictly valid only in the limit \(z\to0\). The difference between
Eqs.~\eqref{eq:15.3.24} and \eqref{eq:15.3.25} vanishes for \(q_0=0\)
and \(q_0=1\), and is less than 10\% for \(0<z<0.5\) (which includes all
galaxies with known red shifts) and \(0<q_0<1.5\). Hence, as long as \(q_0\)
is not very large, there is no substantial difference between using the exact
formula \eqref{eq:15.3.24} and the approximation \eqref{eq:15.3.25} to
determine \(q_0\).

The ``angular diameter distance'' \(d_A\) and the ``proper motion distance''
\(d_M\) are immediately given in terms of \(d_L\) by
Eqs.~\eqref{eq:14.4.22} and \eqref{eq:14.4.23}:
\[
  d_A=(1+z)^{-2}d_L,
  \qquad
  d_M=(1+z)^{-1}d_L,
\]
and the ``parallax distance'' \(d_P\) is given by Eq.~\eqref{eq:14.4.10}
as
\begin{equation}
  d_P
  =
  \frac{
    zq_0+(q_0-1)\bigl(-1+\sqrt{2q_0z+1}\bigr)
  }{
    H_0
    \left[
      q_0^4(1+z)^2
      -(2q_0-1)
      \left\{
        zq_0+(q_0-1)\bigl(-1+\sqrt{2q_0z+1}\bigr)
      \right\}^{\!2}
    \right]^{1/2}
  }.
  \tag{15.3.26}\label{eq:15.3.26}
\end{equation}

The number counts discussed in Section~\ref{sec:14.7} can now be expressed
more explicitly as functionals of the source density \(n\). Shifting variables
from \(t_1\) to \(z\), and using Eqs.~\eqref{eq:15.3.4},
\eqref{eq:15.3.22}, \eqref{eq:15.3.23}, and
\eqref{eq:14.7.7}--\eqref{eq:14.7.9} gives the number of sources with red
shift less than \(z\) and apparent luminosity greater than \(l\) as
\begin{align}
  N(<z,>l)
  ={}&
  \int_0^\infty\dd L
  \int_0^{\min\{z,z_l(L)\}}
  4\pi H_0^{-3}q_0^{-4}
  (1+z')^{-6}(1+2q_0z')^{-1/2}
  \notag\\
  &\quad\times
  \left[
    z'q_0+(q_0-1)\bigl(-1+\sqrt{2q_0z'+1}\bigr)
  \right]^2
  n(z',L)\,\dd z',
  \tag{15.3.27}\label{eq:15.3.27}
\end{align}
where
\begin{equation}
  z_l(L)
  =
  q_0\left(\frac{LH_0^2}{4\pi l}\right)^{1/2}
  +(1-q_0)
  \left[
    -1+
    \left\{
      1+2\left(\frac{LH_0^2}{4\pi l}\right)^{1/2}
    \right\}^{1/2}
  \right],
  \tag{15.3.28}\label{eq:15.3.28}
\end{equation}
and \(n(z,L)\dd L\) is the proper number density of sources at red shift
\(z\) with absolute luminosity between \(L\) and \(L+\dd L\). For radio
sources with the spectrum \eqref{eq:14.7.13},
Eqs.~\eqref{eq:15.3.4}, \eqref{eq:15.3.22}, \eqref{eq:15.3.23},
\eqref{eq:14.7.15}, \eqref{eq:14.7.16}, and \eqref{eq:14.7.8} give the
number of sources with red shift less than \(z\) and intrinsic power of
frequency \(\nu\) greater than \(S\) as
\begin{align}
  N(<z,>S;\nu)
  ={}&
  \int_0^\infty\dd P
  \int_0^{\min\{z,z_{S\alpha}(P)\}}
  4\pi H_0^{-3}q_0^{-4}
  (1+z')^{-6}(1+2q_0z')^{-1/2}
  \notag\\
  &\quad\times
  \left[
    z'q_0+(q_0-1)\bigl(-1+\sqrt{2q_0z'+1}\bigr)
  \right]^2
  n(z',P;\nu)\,\dd z',
  \tag{15.3.29}\label{eq:15.3.29}
\end{align}
where \(z_{S\alpha}(P)\) is the solution of the equation
\begin{equation}
  (1+z)^{(\alpha-1)/2}
  \left[
    zq_0+(q_0-1)\bigl(-1+\sqrt{2q_0z+1}\bigr)
  \right]
  =
  q_0^2H_0\left(\frac{P}{S}\right)^{1/2},
  \tag{15.3.30}\label{eq:15.3.30}
\end{equation}
and \(n(z,P;\nu)\dd P\) is the proper number density of sources at red shift
\(z\) with intrinsic power at frequency \(\nu\) between \(P\) and
\(P+\dd P\). If there were no evolution of sources, then
Eqs.~\eqref{eq:14.7.18} and \eqref{eq:14.7.19} would give \(n\) the
\(z\)-dependence
\[
  n(z,L)=n(0,L)(1+z)^3,
  \qquad
  n(z,P;\nu)=n(0,P;\nu)(1+z)^3,
\]
and the \(z'\) integrals in Eqs.~\eqref{eq:15.3.27} and
\eqref{eq:15.3.29} could be done explicitly. However, we have already seen
in Section~\ref{sec:14.7} that this hypothesis does not agree with
measurements of \(N(>S;\nu)\) for radio sources or \(N(<z)\) for
quasi-stellar sources. Thus Eqs.~\eqref{eq:15.3.27} and
\eqref{eq:15.3.29} can best be used to gain information about the \(z\)
and \(L\) or \(P\) dependence of \(n(z,L)\) or \(n(z,P;\nu)\). In this
way, Longair\hyperref[ch15-ref-15]{\textsuperscript{15}} has concluded that
either the number density of radio sources has been decreasing (apart from the
expansion of the universe) like \(t^{-2.5}\), or the mean source power has
been decreasing like \(t^{-3.5}\). In addition, a sharp cutoff seems to be
needed at early times, although this conclusion is not definitely
established.\hyperref[ch15-ref-16]{\textsuperscript{16}} A study of the
quasi-stellar sources in the 3C catalogue by
Schmidt\hyperref[ch15-ref-17]{\textsuperscript{17}} shows the same general
features---a proper number density that increases with \(z\) much faster than
\((1+z)^3\) for \(0<z<1\), and falls off sharply for \(z>2\). It may be
that this cutoff marks the epoch of galaxy or quasi-stellar source formation.

The age of the universe is one more important datum that can help us to decide
among different cosmological models. A firm lower limit on the age of the
universe is provided by the age of the earth, determined from the relative
abundances of radioactive elements and their decay products in the earth's
crust. In 1929 Lord
Rutherford\hyperref[ch15-ref-18]{\textsuperscript{18}} calculated this age
to be about \(3.4\times10^9\) years. Modern
studies\hyperref[ch15-ref-19]{\textsuperscript{19}} give for the age of the
earth a reliable value of \(4.5\times10^9\) years. If the Hubble time
\(H_0^{-1}\) is \(13\times10^9\) years, then according to
Eq.~\eqref{eq:15.3.11}, the lower limit \(t_0>4.5\times10^9\) years
requires that \(q_0<5\).

Radioactive dating can also be applied to our galaxy. The basic work on the
stellar synthesis of the heavy elements is a 1957 paper by the Burbidges,
Fowler, and Hoyle.\hyperref[ch15-ref-20]{\textsuperscript{20}} (The purpose
of this paper was, at least in part, to defend the steady state model, by
showing that the elements could be formed in stars, without needing a ``big
bang.'' As discussed in Section~\ref{sec:15.7}, it is generally accepted
today that the elements were mostly formed in stars, with the very important
exception that helium may have been formed in the hot, early universe.)
According to this work, the isotopes of uranium were formed by a rapid process
of neutron addition, the \(r\)-process, in an earlier generation of stars.
The initial abundance ratio is
calculated as\hyperref[ch15-ref-21]{\textsuperscript{21}}
\[
  \left[\frac{{}^{235}\mathrm U}{{}^{238}\mathrm U}\right]_1
  =
  1.65\pm0.15
  \qquad\text{(at formation)}.
\]
The decay rates of these isotopes are accurately known to be
\[
  \lambda({}^{235}\mathrm U)
  =0.971\times10^{-9}/\text{year},
  \qquad
  \lambda({}^{238}\mathrm U)
  =0.154\times10^{-9}/\text{year},
\]
and at present their abundance ratio is
\[
  \left[\frac{{}^{235}\mathrm U}{{}^{238}\mathrm U}\right]_0
  =
  0.00723
  \qquad\text{(at present)}.
\]
If all the uranium was formed promptly after the birth of the galaxy at a time
\(t_G\), then the age of the galaxy must
be\hyperref[ch15-ref-20]{\textsuperscript{20}}
\[
  \begin{aligned}
  t_0-t_G
  &=
  \frac{
    \ln[{}^{235}\mathrm U/{}^{238}\mathrm U]_1
    -\ln[{}^{235}\mathrm U/{}^{238}\mathrm U]_0
  }{
    \lambda({}^{235}\mathrm U)-\lambda({}^{238}\mathrm U)
  }\\
  &\approx6.6\times10^9\ \text{years}.
  \end{aligned}
\]

Any society that developed during an earlier epoch in the history of the
galaxy would have found a larger proportion of the fissionable isotope
\({}^{235}\mathrm U\) than now available on earth, and could therefore have
moved toward nuclear destruction even faster than our own civilization.

An error of 20\% in the estimated initial abundance ratio of
\({}^{235}\mathrm U\) and \({}^{238}\mathrm U\) would produce only a 4\%
error in the age of the galaxy. A much greater source of uncertainty arises
from the possibility that appreciable amounts of uranium were formed well
after the galaxy. In this case, the galaxy must be appreciably older than
\(6.6\times10^9\) years. In order to settle this question, other abundance
ratios have been used in conjunction with the ratio of
\({}^{235}\mathrm U\) to \({}^{238}\mathrm U\), the duration of the period
of element synthesis being taken as a free parameter along with the time when
this period began. Using the \({}^{232}\mathrm{Th}/{}^{238}\mathrm U\) and
\({}^{235}\mathrm U/{}^{238}\mathrm U\) ratios, Fowler and
Hoyle\hyperref[ch15-ref-21]{\textsuperscript{21}} estimated that the age of
the oldest \(r\)-process elements is between \(9.6\times10^9\) and
\(15.6\times10^9\) years.
Clayton\hyperref[ch15-ref-22]{\textsuperscript{22}} has included the
\({}^{187}\mathrm{Re}/{}^{187}\mathrm{Os}\) ratio in his analysis, with
results similar to those of Fowler and Hoyle. However, chemical separation
effects may be important here, so these results are subject to possible large
systematic errors. Dicke\hyperref[ch15-ref-23]{\textsuperscript{23}} has
persistently argued that the bulk of the \(r\)-process elements were produced
within a few hundred million years of the formation of the galaxy, in which
case the age of the galaxy would be close to \(7\times10^9\) years. It is
safe to conclude that the galaxy and hence the universe is at least
\(7\times10^9\) years old, so that \(q_0<2.3\) for
\(H_0^{-1}\approx13\times10^9\) years, but radioactive dating cannot yet
be considered to give a precise age for the galaxy.

It is also possible to estimate the age of the galaxy by study of its globular
clusters. These are large compact clusters containing thousands of individual
stars, so their Hertzsprung--Russell diagrams (the relation between luminosity
and spectral type) can be determined with some precision. Also, the low metal
content of the globular cluster stars indicates that they belong to the first
generation of stars to condense out of the protogalaxy (called Population II;
see Section~\ref{sec:14.5}) and are therefore among the oldest objects in our
galaxy. If all stars in a globular cluster have the same initial chemical
composition and age, and differ only in mass, then these stars must fall on a
locus in the Hertzsprung--Russell diagram, whose shape depends only on the age
and initial chemical composition. By comparing computer solutions of the
equations of stellar evolution with the densities of stars in the observed
Hertzsprung--Russell diagrams for a large number of globular clusters,
Iben\hyperref[ch15-ref-24]{\textsuperscript{24}} has deduced cluster ages
ranging from \(8\times10^9\) to \(18\times10^9\) years, corresponding to
initial helium abundances (by mass) ranging from 33 to 24\%. It is not ruled
out that all clusters have the same age, which would most probably lie in the
range from \(9.5\times10^9\) to \(15.5\times10^9\) years. If the age of
the universe is really greater than \(9\times10^9\) years, and if the Hubble
time \(H_0^{-1}\) is \(13\times10^9\) years, then \(q_0\) must be less than
\(\tfrac12\), and the universe must be negatively curved and infinite, as
also indicated by the mass-density estimates discussed in the last section.

It would certainly be premature to reach any definite conclusions about the
curvature of space from these estimates of the age of the universe. However,
the fact that the uranium and globular-cluster ages are roughly comparable with
the Hubble time \(H_0^{-1}\) provides a strong argument that the observed
correlation of red shift with luminosity distance really does have something
to do with the evolution of the universe.

The explicit solution for \(a(t)\) in the matter-dominated era can be used to
illustrate the horizons that limit our vision of the universe. The speed of
light sets an upper limit to the local propagation velocity of any signal, so
at a given time \(t\) an observer at \(r=0\) can receive signals emitted at
time \(t_1\) only from radial coordinates \(r<r_1\), where \(r_1\) is the
radial coordinate from which light signals emitted at time \(t_1\) would just
reach \(r=0\) at time \(t\). According to Eq.~\eqref{eq:14.3.1}, \(r_1\)
is determined by the formula
\begin{equation}
  \int_0^{r_1}\frac{\dd r}{\sqrt{1-kr^2}}
  =
  \int_{t_1}^{t}\frac{\dd t'}{a(t')}.
  \tag{15.3.31}\label{eq:15.3.31}
\end{equation}
If the \(t'\)-integral diverges as \(t_1\to0\), then it is in principle
possible to receive signals emitted at sufficiently early times from any
comoving particle (such as a ``typical galaxy'') in the universe. On the other
hand, if the \(t'\)-integral converges for \(t_1\to0\) (or, in
singularity-free models, for \(t_1\to-\infty\)), then our vision is limited
by what Rindler\hyperref[ch15-ref-25]{\textsuperscript{25}} has called a
\emph{particle horizon}: It is possible to receive signals at time \(t\) only
from comoving particles that lie within the radial coordinate \(r_H(t)\),
where
\[
  \int_0^{r_H(t)}\frac{\dd r}{\sqrt{1-kr^2}}
  =
  \int_0^t\frac{\dd t'}{a(t')}.
\]
The proper distance \eqref{eq:14.2.21} of this horizon is
\begin{equation}
  d_H(t)
  =
  a(t)\int_0^{r_H(t)}\frac{\dd r}{\sqrt{1-kr^2}}
  =
  a(t)\int_0^t\frac{\dd t'}{a(t')}.
  \tag{15.3.32}\label{eq:15.3.32}
\end{equation}
It is easy to see from Eq.~\eqref{eq:15.1.20} that a particle horizon will
be present if \(\rho\) grows faster than \(a^{-2-\epsilon}\) as \(a\to0\),
as would generally be expected. In particular, if the greatest part of the
\(t'\)-integral comes from the matter-dominated era, then
Eq.~\eqref{eq:15.3.4} can be used to express \(\dd t'\) in terms of
\(a(t')\) and \(\dd a(t')\), and we find
\begin{equation}
  d_H(t)
  =
  \begin{cases}
    \displaystyle
    \frac{a(t)}{a_0H_0\sqrt{2q_0-1}}\,
    \cos^{-1}
    \left[
      1-\frac{(2q_0-1)a(t)}{q_0a_0}
    \right],
    &q_0>\tfrac12\quad(k=+1),\\[1.2em]
    \displaystyle
    \frac2{H_0}\left(\frac{a(t)}{a_0}\right)^{3/2},
    &q_0=\tfrac12\quad(k=0),\\[1.2em]
    \displaystyle
    \frac{a(t)}{a_0H_0\sqrt{1-2q_0}}\,
    \cosh^{-1}
    \left[
      1+\frac{(1-2q_0)a(t)}{q_0a_0}
    \right],
    &q_0<\tfrac12\quad(k=-1).
  \end{cases}
  \tag{15.3.33}\label{eq:15.3.33}
\end{equation}
In the early part of the matter-dominated era, \(a\) was much less than
\(a_0\), so the particle horizon was at a small proper distance:
% VERIFY SOURCE: the scan prints t/3. Combining (15.3.4) with the preceding
% expression gives 3t; the printed result is retained for source fidelity.
\begin{equation}
  d_H(t)
  \longrightarrow
  H_0^{-1}
  \left(\frac{q_0}{2}\right)^{-1/2}
  \left(\frac{a}{a_0}\right)^{3/2}
  \approx\frac{t}{3}.
  \tag{15.3.34}\label{eq:15.3.34}
\end{equation}
For \(q_0\le\tfrac12\), \(a(t)\) increases without limit as
\(t\to\infty\), so \(d_H(t)\) increases faster than \(a(t)\), and the
particle horizon will thus eventually expand to include any given comoving
particle. For \(q_0>\tfrac12\), the universe is spatially finite, with a
circumference given by Eq.~\eqref{eq:14.2.4}:
\begin{equation}
  L(t)=2\pi a(t).
  \tag{15.3.35}\label{eq:15.3.35}
\end{equation}
Looking out in any given direction, we can see comoving particles out to a
fraction of this circumference, given by Eqs.~\eqref{eq:15.3.33} and
\eqref{eq:15.2.5} as
\begin{equation}
  \frac{d_H(t)}{L(t)}
  =
  \frac1{2\pi}
  \cos^{-1}
  \left[
    1-\frac{(2q_0-1)a(t)}{q_0a_0}
  \right].
  \tag{15.3.36}\label{eq:15.3.36}
\end{equation}
When \(a(t)\) expands to its maximum value \eqref{eq:15.3.9}, this fraction
will be \(\tfrac12\), and we shall be able to see all the way to the
``antipodes.'' However, this fraction remains less than unity until \(a(t)\)
shrinks once again to zero, so we shall not be able to see all the way around
the universe until then. If \(q_0=1\) and
\(H_0^{-1}=13\times10^9\) years, then the present circumference
\eqref{eq:15.3.35} is given by Eq.~\eqref{eq:15.2.5} as
\(82\times10^9\) light years and the particle horizon is at one-quarter this
distance, or \(20\times10^9\) light years.

Just as there are some comoving particles that we cannot now see, there may in
some cosmological models be events that we never shall see. An event that
occurs at \(r_1\) at time \(t_1\) will become visible at \(r=0\), at a time
\(t\) given by Eq.~\eqref{eq:15.3.31}. If the \(t'\)-integral diverges as
\(t\to\infty\) (or at the time of the next contraction to \(a=0\)), then it
will in principle be possible to receive signals from any event if we wait
long enough. On the other hand, if the \(t'\)-integral converges for large
\(t\), then it will only be possible to receive signals from events for which
\[
  \int_0^{r_1}\frac{\dd r}{\sqrt{1-kr^2}}
  \le
  \int_{t_1}^{t_{\max}}\frac{\dd t'}{a(t')},
\]
where \(t_{\max}\) is either infinity or the time of the next contraction to
\(a=0\). Rindler\hyperref[ch15-ref-25]{\textsuperscript{25}} calls this an
\emph{event horizon}. For \(q_0<\tfrac12\) or \(q_0=\tfrac12\), \(a(t)\)
grows as \(t\to\infty\) like \(t\) or \(t^{2/3}\), so the \(t'\)-integral
diverges at \(t=\infty\), and there is no event horizon. For
\(q_0>\tfrac12\), the \(t'\)-integral converges at \(t_{\max}\), so there
is an event horizon: The only events occurring at time \(t_1\) that will be
visible before the collapse of the universe are those within a proper distance
\begin{align}
  d_E(t_1)
  &=
  a(t_1)\int_{t_1}^{t_{\max}}\frac{\dd t'}{a(t')}
  \notag\\
  &=
  \frac{a(t_1)}{a_0H_0\sqrt{2q_0-1}}
  \left[
    2\pi
    -\cos^{-1}
    \left\{
      1-\frac{(2q_0-1)a(t_1)}{q_0a_0}
    \right\}
  \right].
  \tag{15.3.37}\label{eq:15.3.37}
\end{align}
If \(q_0=1\) and \(H_0^{-1}=13\times10^9\) years, then the only events
occurring now that will ever become visible to us are those that occur within
a proper distance of \(61\times10^9\) light years.
